% TOP_PROBLEM: {"id": 20000313, "title": "A cyclotomic resolution-Betti obstruction for the global canonical index", "category_id": 5, "category": {"id": 5, "name": "algebraic_geometry", "display_name": "Algebraic Geometry", "description": "Geometric objects defined by polynomial equations.", "slug": "algebraic-geometry", "order_index": 5, "created_at": "2026-07-31T15:26:25.671Z"}, "status": "partially_solved"}
% FIRST_POSTED: 2026-10-06
% ATTEMPT_STATUS: unresolved
% OPENAI_CLAIM: {"schema_version":1,"source_repo":"https://github.com/openai/math","source_commit":"adc7f1241b42e322a6451854ab7e4b4c146bf78a","release_date":"2026-10-06","match":"partial","resolution":"partial","independently_reviewed":false,"second_pass":"n/a","families":[{"family":"034","title":"Log abundance for compact Kähler spaces under logarithmic Iitaka subadditivity","subject":"Algebraic and complex geometry","manuscripts":[{"title":"Uniform indices for semi-log-canonical log Calabi–Yau pairs","pdf_path":"preprints/Uniform-indices-for-semi-log-canonical-log-Calabi-Yau-pairs-October-5-2026/uniform-slc-index.pdf","readme_path":"preprints/Uniform-indices-for-semi-log-canonical-log-Calabi-Yau-pairs-October-5-2026/README.md","date":"2026-10-05","theorem_ref":"thm:main (sections/introduction.tex:29)","statement_sources":[{"path":"preprints/Uniform-indices-for-semi-log-canonical-log-Calabi-Yau-pairs-October-5-2026/build/sections/introduction.tex","line":29,"label":"thm:main"}]}],"lean":null,"reasoning_trace":null}]}
\documentclass[11pt]{article}
\usepackage[margin=25mm]{geometry}
\usepackage{fontspec}
\setmainfont{Latin Modern Roman}
\usepackage{amsmath,amssymb,mathtools}
\usepackage{unicode-math}
\setmathfont{Latin Modern Math}
\usepackage{enumitem}
\usepackage{xurl,hyperref,bookmark}
\hypersetup{colorlinks=true,urlcolor=blue}
\setcounter{secnumdepth}{0}
\setlength{\parindent}{0pt}
\setlength{\parskip}{5pt}
\setlength{\emergencystretch}{3em}
\begin{document}
\section*{A cyclotomic resolution-Betti obstruction for the global canonical index}

\subsection{OpenAI's claim}
A dimension- and coefficient-dependent multiple is claimed to trivialize the log canonical line of every connected projective semi-log-canonical Calabi--Yau pair, independently of its number of components. The cited manuscript argues for that result; its relation to this problem is specified below.

\subsection{Mathematical statement of the claimed result}
\paragraph{Mathematical claim 1.} Let $k$ have characteristic zero and be algebraically closed; over $k$, fix $d\ge4$ and a finite $\Phi\subset[0,1]\cap\mathbb Q$. OpenAI asserts an integer $a=a(d,\Phi)>0$ such that every connected projective equidimensional semi-log-canonical pair $(X,B)$ of dimension $d$, with nonzero boundary coefficients in $\Phi$ and $K_X+B\sim_{\mathbb Q}0$, satisfies
\[
a(K_X+B)\text{ is Cartier},\qquad \mathcal O_X(a(K_X+B))\simeq\mathcal O_X.
\]
Semi-log-canonical means $X$ is reduced, $S_2$ and nodal in codimension one, $B$ avoids its nonnormal locus, and the normalization with the conductor boundary is log canonical. The hypothesis is rational linear triviality, not merely numerical triviality.

\subsection{Outline of the argument}
Manuscript 1 (family 034, 2026-10-05): The normal components first receive a uniform pluricanonical exponent. Residues along conductor strata turn descent into a scalar-gluing problem, whose possible characters are bounded using Hodge-norm rank estimates and moving-subvariety fibrations; this produces a common exponent for the glued pair.

\subsection{Scope relative to this problem}
For an algebraically closed characteristic-zero base, the theorem bounds a common trivializing multiple of the adjoint for semi-log-canonical pairs of dimension at least four whose adjoint is already rationally linearly trivial. With zero boundary it covers normal log-canonical varieties in that class. This page assumes numerical triviality, without spelling out a characteristic or a dimension lower bound. Those conventions and the implication from numerical to rational linear triviality must be supplied before claiming full equivalence; the record retains partial coverage.

\subsection{Formal verification}
Family 034 has no Lean scope page in this snapshot; no formalization of these manuscript results is listed.

This site has not executed the Lean files or checked their dependencies. These descriptions report the scope of the linked declarations, rather than a new verification of the proofs.

\subsection{Status caveat}
OpenAI attributes these results to an internal model. Its release notice cautions that results lacking full formalization can contain mistakes. This site has not independently assessed the proofs; a solved label records the scope of the reported claim.

\subsection{Sources}
\begin{itemize}
\item OpenAI, \emph{Uniform indices for semi-log-canonical log Calabi–Yau pairs}, 2026-10-05; citation key \texttt{OAI:Uniform-indices-for-semi-log-canonical-log-Calabi-Yau-pairs-October-5-2026}. \href{https://github.com/openai/math/blob/adc7f1241b42e322a6451854ab7e4b4c146bf78a/preprints/Uniform-indices-for-semi-log-canonical-log-Calabi-Yau-pairs-October-5-2026/uniform-slc-index.pdf}{PDF (GitHub)}; \href{https://github.com/openai/math/raw/adc7f1241b42e322a6451854ab7e4b4c146bf78a/preprints/Uniform-indices-for-semi-log-canonical-log-Calabi-Yau-pairs-October-5-2026/uniform-slc-index.pdf}{PDF (direct)}; \href{https://github.com/openai/math/blob/adc7f1241b42e322a6451854ab7e4b4c146bf78a/preprints/Uniform-indices-for-semi-log-canonical-log-Calabi-Yau-pairs-October-5-2026/README.md}{Manuscript page and citation}; \href{https://github.com/openai/math/blob/main/preprints/Uniform-indices-for-semi-log-canonical-log-Calabi-Yau-pairs-October-5-2026/uniform-slc-index.pdf}{Latest version}; Statement source: \href{https://github.com/openai/math/blob/adc7f1241b42e322a6451854ab7e4b4c146bf78a/preprints/Uniform-indices-for-semi-log-canonical-log-Calabi-Yau-pairs-October-5-2026/build/sections/introduction.tex\#L29}{thm:main}.
\item \href{https://github.com/openai/math/blob/adc7f1241b42e322a6451854ab7e4b4c146bf78a/CONTENTS.md}{OpenAI catalogue entry}.
\end{itemize}
\end{document}
